\label{chap:83-03}

\section{Construcción genealógica del continuo}

\subsection{APP, TRIT y TPK: definiciones rectoras}

La construcción parte de tres objetos irreductibles y causalmente
ordenados. El objeto de \emph{All Possible Paths} (en adelante, \APP)
construye el soporte nonádico común, su grafo de todos los caminos finitos y
las dos hojas de evaluación suma--producto. La unidad trítica (en adelante,
\TRIT) es la unidad local mínima de relación topológica, geométrica,
informacional y, una vez ordenada, temporal: su firma visible posee tres
estados y su levantamiento conserva el acarreo que una lectura binaria o
residual perdería. El \emph{Topological Phase Kernel} (en adelante, \TPK)
es la categoría de transporte que compone esas unidades en rutas y conserva,
dentro del mismo estado, hoja, orientación, frontera, registro genealógico, memoria, fase y
supervivencia.

\begin{center}
\begin{tabularx}{.94\textwidth}{@{}p{.10\textwidth}Y Y@{}}
\toprule
Objeto & Construcción & Invariantes transportados\\
\midrule
APP & carta positiva, toro finito, grafo de Cayley, caminos y hojas
\(\Sigma,\Pi\) & marca, celda, vecindad, suma, producto, cociente y cierre\\
TRIT & firma ternaria visible y levantamiento topológico--temporal local &
residuo, carry, hoja, orientación, umbral y reversión\\
TPK & categoría de rutas, fibras y extensiones compatibles & cilindro,
frontera, registro genealógico, memoria, fase, holonomía y supervivencia\\
\bottomrule
\end{tabularx}
\end{center}

La construcción exacta de APP se establece sobre un segmento marcado y una
carta finita; la recta real comparece después como publicación límite. El
intervalo de referencia, sus marcas y el camino orientado que las recorre son:
\begin{equation}
I_{10}=[0,10],\qquad
\mathcal M_{10}=I_{10}\cap\mathbb Z=\{0,1,\ldots,10\},
\qquad
\mathsf P_{10}:0\longrightarrow1\longrightarrow\cdots
\longrightarrow9\longrightarrow10,
\label{p12-83-c3-s1:eq:app-marked-segment}
\end{equation}
Los extremos fijan el borde y las nueve marcas interiores
\(I_{10}^{\circ}\cap\mathbb Z\) forman el alfabeto positivo. El objeto
matemático exacto de partida es la carta
\begin{equation}
\mathcal D_9=\{1,\ldots,9\},
\qquad \mathcal D_8=\{1,\ldots,8\},
\qquad
\lambda_9:\mathcal D_9\xrightarrow{\sim}\mathbb Z/9\mathbb Z,
\qquad \lambda_9(9)=[0]_9.
\label{p12-83-c3-s1:eq:app-positive-chart}
\end{equation}
La marca \(9\) es simultáneamente el último representante positivo y la
publicación de la clase nula. Para todo \(n\geq1\), raíz digital positiva y
cociente se separan mediante
\begin{equation}
\operatorname{dr}_9(n)=1+((n-1)\bmod9),
\qquad q_9(n)=\left\lfloor\frac{n-1}{9}\right\rfloor,
\qquad n=9q_9(n)+\operatorname{dr}_9(n).
\label{p12-83-c3-s1:eq:positive-digital-root-carry}
\end{equation}
Al camino \(1\to2\to\cdots\to9\) se añade la
única arista de cierre \(9\to1\). El sucesor y su memoria se separan:
\begin{equation}
u_+(a)=\begin{cases}a+1,&a<9,\\1,&a=9,\end{cases}
\qquad
\kappa_+(a)=\begin{cases}0,&a<9,\\1,&a=9,\end{cases}
\qquad a+1=u_+(a)+9\kappa_+(a).
\label{p12-83-c3-s1:eq:app-successor-carry}
\end{equation}
Por ello \(9\to1\) cierra la fase, aumenta la memoria y prolonga el estado
genealógico.

El cuadrado de la carta produce el toro finito y su grafo de movimientos:
\begin{equation}
X_9=(\mathbb Z/9\mathbb Z)^2,
\qquad
\Gamma_{\rm APP}=\operatorname{Cay}
\bigl(X_9,\{\pm e_1,\pm e_2\}\bigr),
\qquad
\operatorname{Path}(\Gamma_{\rm APP}).
\label{p12-83-c3-s1:eq:app-cayley-torus}
\end{equation}
La identificación de lados opuestos mantiene el cardinal de la carta y añade
la ley de continuación. El núcleo de \(8^2=64\) intersticios se completa con el gnomon
mínimo que publica la clase nula en las dos coordenadas:
\begin{equation}
\mathcal D_9^2=\mathcal D_8^2\sqcup\mathcal G_9,
\qquad
81=64+(9+8)=64+17.
\label{p12-83-c3-s1:eq:app-gnomon-public}
\end{equation}
La misma ley se prolonga a cartas de mayor profundidad; la escala cambia,
pero el cierre, el borde compartido y la memoria de vuelta conservan su tipo.

La completación volumétrica asociada se expresa mediante la identidad de
estratos
\begin{equation}
10^3=(9+1)^3=729+243+27+1.
\label{p12-83-c3-s1:eq:completion-unit}
\end{equation}
\input{figures/fig_app_emrh_mini}
El volumen interior, las caras, las aristas y el cierre puntual quedan así
tipados como capas distintas de una misma unidad completada. Sobre las 81
celdas de la carta positiva actúan las evaluaciones
\[
\Sigma,\Pi:\mathcal D_9^2\longrightarrow\mathcal D_9,
\qquad
\Sigma(i,j)=\operatorname{dr}_9(i+j),
\qquad
\Pi(i,j)=\operatorname{dr}_9(ij),
\]
Sus residuos y cocientes levantados se definen por
\[
\begin{aligned}
R^+(i,j)&:=\Sigma(i,j),&
Q^+(i,j)&:=\frac{i+j-R^+(i,j)}9,\\
R^\times(i,j)&:=\Pi(i,j),&
Q^\times(i,j)&:=\frac{ij-R^\times(i,j)}9,
\end{aligned}
\]
y, por tanto,
\[
\begin{aligned}
\widehat\Sigma(i,j)&:=\bigl(R^+(i,j),Q^+(i,j)\bigr),\\
\widehat\Pi(i,j)&:=\bigl(R^\times(i,j),Q^\times(i,j)\bigr),
\end{aligned}
\qquad
\widehat\Sigma,\widehat\Pi:
\mathcal D_9^2\longrightarrow\mathcal D_9\times\mathbb Z_{\geq0}.
\]
La biyección \(\lambda_9^{\times2}:\mathcal D_9^2\to X_9\) transporta
estas operaciones y sus levantamientos al toro. Por tanto, APP es la tupla
\begin{equation}
\APP=\bigl(\mathcal D_9,\lambda_9,X_9,\Gamma_{\rm APP},
\operatorname{Path}(\Gamma_{\rm APP}),
\Sigma,\Pi,\widehat\Sigma,\widehat\Pi\bigr),
\label{p12-83-c3-s1:eq:app-full-definition-public}
\end{equation}
La tabla constituye una representación bidimensional de esta tupla
completa. La hoja aditiva es latina; la multiplicativa conserva
el radical ternario del anillo \(\mathbb Z/9\mathbb Z\)
\begin{equation}
R_9:=\mathbb Z/9\mathbb Z,
\qquad
\mathfrak m:=([3]_9)=\{[0]_9,[3]_9,[6]_9\}\triangleleft R_9,
\qquad \mathfrak m^2=(0).
\label{p12-83-c3-s1:eq:app-radical-369}
\end{equation}
La rotación interna es el automorfismo
\[
\rho:\mathfrak m^3\longrightarrow\mathfrak m^3,
\qquad \rho(x_1,x_2,x_3)=(x_2,x_3,x_1),
\]
cuya órbita distinguida es
\[
([3]_9,[6]_9,[0]_9)\xrightarrow{\rho}
([6]_9,[0]_9,[3]_9)\xrightarrow{\rho}
([0]_9,[3]_9,[6]_9)\xrightarrow{\rho}
([3]_9,[6]_9,[0]_9).
\]
La reducción directa envía \(3,6,9\) a \(0,0,0\). La coordenada interna
del ideal conserva, en cambio, el factor extraído:
\begin{equation}
\eta_{\mathfrak m}:(\mathfrak m,+)
\xrightarrow{\;\sim\;}(\mathbb Z/3\mathbb Z,+),
\qquad
\eta_{\mathfrak m}([3a]_9):=[a]_3,
\qquad
([3]_9,[6]_9,[0]_9)
\xmapsto{\;\eta_{\mathfrak m}^{\times3}\;}
([1]_3,[2]_3,[0]_3),
\label{p12-83-c3-s1:eq:app-369-120}
\end{equation}
de modo que el sector \(3,6,9\) permanece íntegro y su coordenada interna se
lee como \(1,2,0\). Esta aplicación pertenece al radical de APP; la
proyección temporal de TRIT constituye el lector calendárico posterior. Las
tres clases calendáricas
quedan separadas como
\[
\mathcal C_+=\{1,4,7\},\qquad
\mathcal C_-=\{2,5,8\},\qquad
\mathcal C_0=\{3,6,9\}.
\]
La incompatibilidad entre ambas hojas se vuelve un invariante medible. Sobre
\(\mathscr A_3=\mathcal D_9^3\), con medida uniforme
\(\nu_3(x)=9^{-3}\), sea
\(\mathcal H_3=L^2(\mathscr A_3,\nu_3)\). Definimos
\[
\Sigma_3,\Pi_3:\mathscr A_3\longrightarrow\mathcal D_9,
\qquad
\Sigma_3(x)=\operatorname{dr}_9(x_1+x_2+x_3),
\qquad
\Pi_3(x)=\operatorname{dr}_9(x_1x_2x_3),
\]
y los proyectores ortogonales
\[
P_{\Sigma,3}:=\mathbb E_{\nu_3}[\,\cdot\mid\sigma(\Sigma_3)],
\qquad
P_{\Pi,3}:=\mathbb E_{\nu_3}[\,\cdot\mid\sigma(\Pi_3)]
\in\mathcal B(\mathcal H_3).
\]
La enumeración conjunta de las \(9^3\) palabras determina entonces
\[
\bigl\|[P_{\Sigma,3},P_{\Pi,3}]\bigr\|=\frac{\sqrt3}{4},
\qquad
B_c=7I+2R=9P_++5P_-,
\qquad
459-405=54.
\]
El primer número reaparece como prefactor electrónico; el segundo organiza
el bloque central; el tercero registra el defecto global suma--producto.

\subsection{TRIT: unidad local topológica y unidad de relación temporal}

HMT fija primero la firma visible
\begin{equation}
\mathrm{TRIT}_{\rm vis}=\{-1,0,+1\},
\qquad \tau\in\mathrm{TRIT}_{\rm vis}.
\label{p12-83-c3-s1:eq:trit-visible-signature}
\end{equation}
Frente a la unidad binaria, el TRIT conserva tres regímenes locales. En la
distribución uniforme, su información y su entropía se distinguen como
\begin{equation}
I_{\rm TRIT}=\log3\ \text{nats}=\log_2 3\ \text{bits},
\qquad S_{\rm TRIT}=k_BI_{\rm TRIT}.
\label{p12-83-c3-s1:eq:trit-information-unit}
\end{equation}
El bit sigue siendo la unidad binaria; el TRIT es la unidad mínima capaz de
distinguir simultáneamente retorno, umbral y apertura. La igualdad
\eqref{p12-83-c3-s1:eq:trit-information-unit} cuantifica la información de la terna; el
coste \(k_B\Theta\log3\) pertenece al lector térmico de reinicio posterior.
Mecánica Dimensional conserva además el antecedente aritmético de esa firma.
Para \((i,j)\in\mathcal D_9^2\), la celda aritmética levantada asociada es
\begin{equation}
c_{\rm MD}(i,j;\tau)=
\bigl(i,j;R^+,Q^+,R^\times,Q^\times;\tau\bigr).
\label{p12-83-c3-s1:eq:trit-complete-local-unit}
\end{equation}
Los residuos \(R^+,R^\times\) son la publicación visible de las dos hojas;
los cocientes \(Q^+,Q^\times\) conservan la memoria exacta de la división
euclídea; y \(\tau\) orienta la transición. La firma visible ocupa el primer
nivel; la celda que la porta añade los antecedentes aritméticos, y el estado
enriquecido completa después frontera, ruta, fase y supervivencia.

Conviene separar cuatro estratos. La firma visible es \(\tau\); el par
residuo--acarreo es su levantamiento local; \(c_{\rm MD}\) es la celda
aritmética de APP provista de esa firma; y el estado TPK añade extensión,
frontera, registro genealógico, fase y supervivencia. Esta jerarquía conserva la identidad
propia de la coordenada observada, su levantamiento, la celda portadora y la
historia enriquecida que la transporta.

En la carta ternaria balanceada, el entero local levantado
\(n\in\{-4,-3,\ldots,3,4\}\) se descompone de manera única como
\begin{equation}
\widehat n=(r,\kappa),
\qquad n=r+3\kappa,
\qquad r,\kappa\in\{-1,0,+1\}.
\label{p12-83-c3-s1:eq:trit-balanced-lift}
\end{equation}
La firma visible es \(\tau=r\); \(\kappa\) conserva el carry que la proyección
\(n\mapsto r\) olvidaría.
La fibra sobre el cero visible contiene tres estados distintos,
\[
0^+=(0,+1),\qquad 0^0=(0,0),\qquad 0^-=(0,-1),
\]
porque \(n=-3,0,3\) poseen el mismo residuo visible y carries distintos. Éste
es el levantamiento topológico mínimo de MD: una transición orientada de APP
queda provista de la información necesaria para invertirla y para decidir qué
prolongaciones siguen siendo compatibles.

La firma posee además una regla de memoria de hoja. Si \(m_k\) es el modo
aritmético activo, entonces
\begin{equation}
m_{k+1}=
\begin{cases}
\Sigma,&\tau_k=+1,\\
\Pi,&\tau_k=-1,\\
m_k,&\tau_k=0.
\end{cases}
\label{p12-83-c3-s1:eq:trit-zero-memory}
\end{equation}
El cero codifica retención de la hoja anterior. Por eso la palabra
\((+1,-1,0)\) publica la secuencia \((\Sigma,\Pi,\Pi)\): el tercer estado
conserva el modo multiplicativo activo.

La relación temporal se obtiene al parametrizar ordenadamente estas celdas.
Cada firma porta un álgebra cuadrática
\begin{equation}
\mathbb A_\tau=\mathbb R[X]/(X^2+\tau),
\qquad J_\tau=[X],
\label{p12-83-c3-s1:eq:trit-quadratic-algebra-local}
\end{equation}
de modo que
\[
\begin{array}{c@{\qquad}c@{\qquad}l}
\toprule
\tau & J_\tau^2 & \text{relación temporal ordenada}\\
\midrule
+1 & -I & \text{retorno elíptico, memoria y pasado},\\
0  & 0  & \text{umbral parabólico y presente},\\
-1 & +I & \text{apertura hiperbólica bidireccional y futuro}.\\
\bottomrule
\end{array}
\]
El TRIT precede al tiempo exterior: primero clasifica el tipo algebraico y
topológico del transporte; después, una parametrización ordenada lo publica
como relación temporal. La doble dirección queda tipada por la involución
exacta
\begin{equation}
\mathcal C_{\rm rev}(\tau_1,\ldots,\tau_n)
=(-\tau_n,\ldots,-\tau_1),
\qquad \mathcal C_{\rm rev}^2=I,
\label{p12-83-c3-s1:eq:trit-time-reversal}
\end{equation}
que intercambia retorno y apertura y conserva el umbral. En la hoja
hiperbólica, \(P_\pm=(I\pm J_{-1})/2\) separan las dos direcciones propias y
la reversión intercambia sus orientaciones. Topología y tiempo son, así, las
dos lecturas coordinadas del TRIT.

La involución combinatoria \(\mathcal C_{\rm rev}\) y una conjugación
operatorial \(C\) pertenecen, sin embargo, a dominios distintos. En una
realización dinámica separada, provista de \(H=H^*\),
\(CJ_EC^{-1}=-J_E\) y \(CHC^{-1}=H\), fijamos una escala de acción interna
\(\hbar_{\rm int}>0\) y escribimos
\(h_{\rm int}:=2\pi\hbar_{\rm int}\). Entonces
\begin{equation}
U(t)=\exp(-J_EHt/\hbar_{\rm int}),
\qquad CU(t)C^{-1}=U(-t).
\label{p12-83-c3-s1:eq:trit-bidirectional-propagator}
\end{equation}
El entrelazador que transporta palabras TRIT a la realización dinámica fija
la correspondencia entre \(\mathcal C_{\rm rev}\) y \(C\). La doble
orientación local proporciona el dominio de esa construcción.

\subsection{TPK: categoría de transporte y evolución de fase}

El \emph{Topological Phase Kernel} es la categoría de transporte sobre la
base observable de APP que compone cifras, palabras, operadores y relaciones
de prolongación en un mismo estado enriquecido. Su arquitectura interna se distribuye en
ocho clases operatorias, diferenciadas por dominio, codominio y función:
\begin{enumerate}[label=\textup{(\roman*)},leftmargin=2.25em]
\item soporte observable y estado enriquecido;
\item selector interno de hoja \(M_{\rm ph}:\{1,\ldots,9\}\to\{+1,-1,0\}\);
\item traslaciones cardinales y evolución visible \(F_{\rm obs}\);
\item lectores y cocientes nonádico, trítico y de bloque \(1000\);
\item memoria, acarreo, cilindro, frontera y registro genealógico;
\item periodicidad y retorno
\(C_{12}\hookrightarrow C_{108}\twoheadrightarrow C_9\),
\(\mathcal K_{27}\) y canales \(90/120\);
\item prolongaciones \(\operatorname{Ext}_r\), cilindros, supervivencia y
extensión natural bidireccional;
\item funtores de salida arquimediana, excepcional calibrada y dimensional.
\end{enumerate}
El retorno de 27 pasos es el morfismo compuesto
\begin{equation}
\mathcal K_{27}:=F_{\rm obs}^{27},
\qquad \mathcal K_{27}^{4}=F_{\rm obs}^{108}.
\label{p12-83-c3-s1:eq:tpk-kernel-27}
\end{equation}
Su cuarto iterado restaura la fase observable del calendario, mientras la
memoria y el estado TPK completo continúan su trayectoria;
\(\mathcal K_{27}\) es el morfismo de retorno interno de \(C_{108}\).
El calendario \(C_{108}\), la conexión nonádica, la extensión bilateral y
los lectores ocupan lugares tipados dentro de esta arquitectura. Sea
\(\mathsf P_{\rm APP}\) la categoría
libre del grafo dirigido de APP y
\(\mathsf P_{\rm APP}^{(2)}=\mathsf P_{\rm APP}\times
\mathsf P_{\rm APP}\) su lectura simultánea suma--producto. Sea además
\[
\mathcal Q_{\rm obs}
=X_9^2\times\{N,E,S,O\}^2\times\Theta_{108},
\qquad \Theta_{108}=\mathbb Z/108\mathbb Z,
\]
donde, para el representante
\(\widetilde\theta\in\{0,\ldots,107\}\),
\[
r(\theta)=1+(\widetilde\theta\bmod9),
\qquad g_0(\theta)=[\widetilde\theta]_9,
\]
\[
\tau(\theta)=
\begin{cases}
+1,&r(\theta)\in\{1,4,7\},\\
-1,&r(\theta)\in\{2,5,8\},\\
0,&r(\theta)\in\{3,6,9\},
\end{cases}
\qquad
\mu(\theta)=
\begin{cases}
\Sigma,&\tau(\theta)=+1,\\
\Pi,&\tau(\theta)=-1,\\
\varnothing,&\tau(\theta)=0.
\end{cases}
\]
La etiqueta calendárica \(\mu=\varnothing\) señala el umbral, y la memoria
dinámica \(m_k\) de \eqref{p12-83-c3-s1:eq:trit-zero-memory} retiene la última hoja
activa. Una configuración
\(q=(p^\Sigma,p^\Pi,d^\Sigma,d^\Pi,\theta)\) conserva las dos posiciones,
las dos direcciones cardinales y el índice de calendario. La evolución
visible \(F_{\rm obs}\) activa la hoja indicada por TRIT, transporta la
posición correspondiente y avanza \(\theta\). Denotemos por
\(\mathsf P_{\rm obs}\) la categoría libre generada por las flechas
\(q\to F_{\rm obs}(q)\). Cada flecha induce un par de caminos de APP y
define el funtor de base
\begin{equation}
B:\mathsf P_{\rm obs}\longrightarrow\mathsf P_{\rm APP}^{(2)}.
\label{p12-83-c3-s1:eq:tpk-bivariant-base-functor}
\end{equation}

Sobre cada \(q\) se dispone la fibra tipada
\begin{equation}
\mathcal F_q=
\mathsf O_q\times\mathsf H_q\times\mathsf C_q\times\mathsf S_q
\times\mathsf B_q\times\mathsf\Lambda_q,
\label{p12-83-c3-s1:eq:tpk-fiber}
\end{equation}
Sus seis factores conservan, respectivamente, orientación y hoja; residuo y
cociente; acarreo; firma; cilindro y frontera; y registro ordenado de
transporte. Memoria y supervivencia son propiedades de la composición y de
las relaciones de prolongación.

Para una ruta \(r:q\to q'\), la relación
\begin{equation}
\operatorname{Ext}_r\subseteq\mathcal F_q\times\mathcal F_{q'}
\label{p12-83-c3-s1:eq:tpk-extension-relation}
\end{equation}
contiene todas las continuaciones admisibles y satisface
\[
\Delta_q\subseteq\operatorname{Ext}_{\operatorname{id}_q},
\qquad
\operatorname{Ext}_{r_2}\circ\operatorname{Ext}_{r_1}
\subseteq\operatorname{Ext}_{r_2\circ r_1}.
\]
Los transportes respetan identidades y composición:
\begin{equation}
\begin{gathered}
\mathsf H(\operatorname{id}_q)=I,
\quad \mathsf C(\operatorname{id}_q)=I,
\quad \mathsf\Lambda(\operatorname{id}_q)=I,
\quad \kappa(\operatorname{id}_q)=0,\\
\mathsf H(r_2\!\circ r_1)=\mathsf H(r_2)\mathsf H(r_1),
\quad
\mathsf C(r_2\!\circ r_1)=\mathsf C(r_2)\mathsf C(r_1),
\quad
\mathsf\Lambda(r_2\!\circ r_1)=
\mathsf\Lambda(r_2)\mathsf\Lambda(r_1).
\end{gathered}
\label{p12-83-c3-s1:eq:tpk-functorial-transport}
\end{equation}
Junto con estas leyes functoriales, las coordenadas transportadas obedecen
\begin{equation}
\begin{aligned}
h'&=\mathsf H(r)(h),\\
c'&=\mathsf C(r)(c)+\kappa(r),\\
\lambda'&=\mathsf\Lambda(r)(\lambda),\\
\sigma'&=\operatorname{Sig}_{q'}(h',c',\lambda'),
\end{aligned}
\qquad
\kappa(r_2\circ r_1)
=\kappa(r_2)+\mathsf C(r_2)\kappa(r_1).
\label{p12-83-c3-s1:eq:tpk-typed-transport}
\end{equation}
La última identidad es el cociclo de acarreo que hace asociativa la
composición con memoria.

El TPK es entonces la categoría \(\mathsf{TPK}\) cuyos objetos son
\begin{equation}
x=(q,o,h,c,\sigma,C,b,\lambda),
\qquad (o,h,c,\sigma,(C,b),\lambda)\in\mathcal F_q,
\label{p12-83-c3-s1:eq:tpk-enriched-object}
\end{equation}
y cuyos morfismos \(x\to x'\) son las ternas \((r,x,x')\) tales que
las coordenadas satisfacen \eqref{p12-83-c3-s1:eq:tpk-typed-transport} y
\((f_x,f_{x'})\in\operatorname{Ext}_r\), donde
\(f_x\in\mathcal F_q\) y \(f_{x'}\in\mathcal F_{q'}\). La composición es
\begin{equation}
(r_2,x',x'')\circ(r_1,x,x')
=(r_2\circ r_1,x,x''),
\label{p12-83-c3-s1:eq:tpk-morphism-composition}
\end{equation}
\Needspace{7\baselineskip}
la identidad es \((\operatorname{id}_q,x,x)\), y el funtor
\begin{equation}
p:\mathsf{TPK}\longrightarrow\mathsf P_{\rm obs},
\qquad p(x)=q,
\qquad p(r,x,x')=r,
\label{p12-83-c3-s1:eq:tpk-base-functor}
\end{equation}
publica la configuración visible sin identificar estados de la misma fibra.
La definición fija una categoría sobre una base; no presupone que sea una
fibración categórica ni que toda flecha admita un levantamiento cartesiano.
Las proyecciones locales quedan separadas por tipo:
\[
\pi_{\rm TRIT}:\operatorname{Ob}(\mathsf{TPK})
\longrightarrow\mathrm{TRIT}_{\rm vis},
\qquad
\pi_{\rm TRIT}(x)=\tau(\theta),
\]
\[
\pi_{\rm loc}(x)=
\bigl(\pi_{\rm TRIT}(x),\mu(\theta),o,h,c,
\sigma,\lambda,g_0(\theta)\bigr).
\]
La primera recupera únicamente la firma TRIT visible. La segunda conserva
además la etiqueta calendárica de hoja, la orientación, el par
residuo--cociente, el acarreo, la firma de transporte, el registro genealógico y la fase,
pero todavía olvida el cilindro, su etiqueta de borde y otras coordenadas de
\(q\). Ésta es la razón por la que ni el TRIT visible ni su registro local
sustituyen al estado enriquecido del TPK.

El nombre recoge tres funciones inseparables. \emph{Topological} remite a
la categoría de caminos, la incidencia, el borde y la prolongación por
cilindros; \emph{Phase} remite al índice nonádico, al calendario orientado y
a la holonomía; \emph{Kernel} nombra el núcleo persistente de transporte que
permanece antes de contraer el estado mediante un lector. Formalmente, el
objeto es la categoría anterior: una topología adicional sobre objetos o
morfismos, o su identificación con el núcleo lineal de un mapa, requiere
datos independientes.

En consecuencia, dos rutas con el mismo extremo visible pueden seguir siendo
objetos distintos por hoja, cociente, acarreo, frontera, cilindro, fase o
registro. El TPK no elige una salida escalar: conserva el antecedente común
del que parten las publicaciones arquimediana, excepcional y dimensional.

% Estas consecuencias se publican después de la construcción correlativa y
% del paso al límite de los capítulos 4 y 5. Se definen aquí para preservar
% íntegramente su procedencia, pero no se ejecutan antes de que exista el
% continuo al que se aplican.
\long\def\HMTDeferredContinuumConsequences{%
\section{Consecuencias y publicaciones del continuo construido}

\subsection{Teorema de ensamblaje terminal de las construcciones correlativas generadas por el estado común}

Cada profundidad \(k\) posee un estado enriquecido \(\mathcal E_k\) con
todas sus emisiones supervivientes y un único constructor correlativo
\begin{equation}
\mathcal O_k:\mathcal E_k\longrightarrow\mathcal Z_k,
\qquad
\widehat x_k\longmapsto
\mathcal O\bigl(\{\mathfrak e_\alpha:
\alpha\in\operatorname{Ch}_k(\widehat x_k)\}\bigr).
\label{p12-83-c3-s1:eq:five-features-single-constructor}
\end{equation}
El valor de \(\mathcal O_k\) conserva simultáneamente las coordenadas
espectral, solenoidal, gauge, coherente y determinantal,
\begin{equation}
(\mathrm{sp},\mathrm{sol},\mathrm{gau},\mathrm{coh},\mathrm{det}),
\label{p12-83-c3-s1:eq:five-features-coordinates}
\end{equation}
como cinco características del mismo estado y del mismo registro genealógico. Los
truncamientos
\[
u^{\mathfrak e}_{\ell,k}:\mathcal E_\ell\to\mathcal E_k,
\qquad
u^{\mathcal O}_{\ell,k}:\mathcal Z_\ell\to\mathcal Z_k
\]
satisfacen la naturalidad única
\begin{equation}
u^{\mathcal O}_{\ell,k}\circ\mathcal O_\ell
=\mathcal O_k\circ u^{\mathfrak e}_{\ell,k},
\qquad \ell\ge k.
\label{p12-83-c3-s1:eq:five-features-naturality}
\end{equation}
Por ello, la acción correlativa pasa al límite como
\begin{equation}
\mathcal O_\infty:
\mathcal C_{\rm cont}^{\rm disc}:=\varprojlim_k\mathcal E_k
\longrightarrow
\mathcal Z_\infty:=\varprojlim_k\mathcal Z_k.
\label{p12-83-c3-s1:eq:five-features-infinite-operator}
\end{equation}
Las cinco aplicaciones terminales son vistas intrínsecas de ese único valor:
\begin{equation}
\boxed{
\mathfrak G_T^{\rm int}:=
\operatorname{View}_T\!\left(
\mathcal O_\infty(\mathcal C_{\rm cont}^{\rm disc})\right),
\quad
T\in\{\mathrm{sp},\mathrm{sol},\mathrm{gau},
\mathrm{coh},\mathrm{det}\}.}
\label{p12-83-c3-s1:eq:five-features-views}
\end{equation}
La simultaneidad precede a las resoluciones especializadas: cada vista expone
la coordenada requerida y el estado completo conserva las otras cuatro.

\subsection{Supervivencia y límite inverso}

Sea \(\operatorname{Surv}^{(\chi)}_N\) el conjunto de historias de
profundidad \(N\) que satisfacen simultáneamente las leyes activas del canal
\(\chi\) y admiten prolongación. Definimos las truncaciones compatibles
\[
\tau_{N+1,N}:\operatorname{Surv}^{(\chi)}_{N+1}
\longrightarrow\operatorname{Surv}^{(\chi)}_N
\]
Estas conservan el prefijo enriquecido. El continuo genealógico es el límite
\begin{equation}
X_\chi=\varprojlim_N
\bigl(\operatorname{Surv}^{(\chi)}_N,\tau_{N+1,N}\bigr),
\qquad
\tau_{N+1,N}(x_{N+1})=x_N.
\label{p12-83-c3-s1:eq:inverse-survival}
\end{equation}
El espacio global de secciones supervivientes y una de sus secciones
enriquecidas se denotan, respectivamente, por
\begin{equation}
X_\infty:=\bigsqcup_\chi X_\chi,
\qquad
\widehat\Omega_\infty\in X_\infty.
\label{p12-83-c3-s1:eq:global-survival-space}
\end{equation}
El lector arquimediano global es
\begin{equation}
P_{\rm arq}:=\operatorname{ev}_{\mathbb R}:
X_\infty\longrightarrow\mathbb R;
\label{p12-83-c3-s1:eq:archimedean-reader}
\end{equation}
su restricción a cada canal evalúa la intersección unitaria de los cilindros
compatibles de la sección correspondiente.
Cuando interviene el lector excepcional, \(X_\infty^{\rm cal}\subseteq
X_\infty\) denota el subespacio en el que se ha fijado el tipo de calibración
de \eqref{p12-83-c3-s1:eq:exceptional-calibration}.
La supervivencia determina la sección por compatibilidad interna: cada
historia debe ser prolongable en el sistema completo. El lector numérico
interviene después para certificar y publicar la intersección del
descendiente ya compatible.

\subsection{Conexión nonádica conjunta, holonomía y memoria vertical}

En profundidad \(K\), la holonomía del retorno nonádico es inicialmente la
correspondencia tipada
\begin{equation}
\Gamma_{9,K}=\operatorname{Hol}_{C_9}(\nabla_{\rm TPK}):
\widehat{\mathcal X}_K\rightrightarrows
\widehat{\mathcal X}_{K+9}.
\label{p12-83-c3-s1:eq:gamma9}
\end{equation}
Sus iterados son composiciones relacionales. Sólo cuando cada estado posee
una prolongación compatible única se restringe a un mapa, y sólo cuando ese
mapa es invertible constituye una monodromía. Esta distinción conserva todas
las extensiones supervivientes antes de seleccionar una sección.
Tras nueve fases retorna la fase local, mientras el estado completo avanza:
\[
g(K+9)=g(K),\qquad q_{{\rm mem},K+9}=q_{{\rm mem},K}+1,
\qquad \widehat x_{K+9}\ne\widehat x_K\ \text{en general}.
\]
Bloque, cilindro, prefijo, acarreo y memoria continúan evolucionando. Ésta
es la diferencia entre periodicidad visible y retorno con biografía.

El estado transportado en el nivel (K) puede exhibirse como
\begin{equation}
\widehat x_K=
(B_K,C_K,p_K,c_K,\Sigma_K,W_K,g(K),q_{{\rm mem},K},s_K),
\label{p12-83-c3-s1:eq:enriched-state-k}
\end{equation}
donde cada coordenada tiene una ley de transporte propia. El retorno de fase
se representa sobre el catálogo
\[
\mathcal U_6^{\rm cat}\cong\mathfrak S_+\times\mathfrak S_\times,
\qquad |\mathfrak S_+|=18,
\qquad |\mathfrak S_\times|=26.
\]
Los promedios condicionales
\[
(E_+f)(s,e)=\frac1{18}\sum_{s'\in\mathfrak S_+}f(s',e),
\qquad
(E_\times f)(s,e)=\frac1{26}\sum_{e'\in\mathfrak S_\times}f(s,e')
\]
conmutan. La palabra de fase
\(\tau=(+1,-1,0,+1,-1,0,+1,-1,0)\) aplica
\(E_+,E_\times,I\) en ese orden, y una vuelta se cierra mediante
\begin{equation}
\mathcal K_{\rm ph}^{\,9}
=(E_+E_\times)\otimes I_{\ell^2(C_9)},
\qquad
\bigl\langle\delta_u,E_+E_\times\delta_v\bigr\rangle
=\frac1{468}\quad(u,v\in\mathcal U_6^{\rm cat}).
\label{p12-83-c3-s1:eq:joint-nine-holonomy}
\end{equation}
Las hojas suma y producto no actúan como filtros independientes: son dos
componentes de una única holonomía que sólo se cierra después de recorrer el
ciclo completo.

La coligación compresión--expresión separa lo que vuelve y lo que permanece
como memoria. Si (J) incrusta el espacio de historias en el espacio de
fase, (T) es la compresión superviviente y \(\eta\) el canal de defecto,
entonces
\begin{equation}
\boxed{
U_{\Gamma_9}J=JT+\eta,
\qquad J^*\eta=0,
\qquad I-T^*T=\eta^*\eta.}
\label{p12-83-c3-s1:eq:compression-expression}
\end{equation}
La primera igualdad descompone el transporte en retorno comprimido y
expresión de memoria; la segunda prueba su ortogonalidad; la tercera mide
exactamente la información que no cabe en la sombra periódica. Así queda
formalizada la afirmación central: vuelve la fase, no el estado completo.

La memoria vertical se organiza mediante la extensión cíclica orientada
\begin{equation}
0\longrightarrow C_{12}\xrightarrow{\,\iota\,}C_{108}
\xrightarrow{\,p\,}C_9\longrightarrow0,
\qquad
c(r,r')=\left[\left\lfloor\frac{r+r'}9\right\rfloor\right]_{12}.
\label{p12-83-c3-s1:eq:c108-extension}
\end{equation}
Aquí \(\iota([b]_{12})=[9b]_{108}\) y
\(p([\theta]_{108})=[\theta]_9\). La extensión no se escinde:
\(C_{108}\) tiene exponente 108, mientras
\(C_{12}\times C_9\) tiene exponente 36.
El cociclo registra el acarreo de cada vuelta de la base \(C_9\) en la
memoria dodecafásica \(C_{12}\). \(C_{108}\) es, por tanto, un calendario
orientado de fase y memoria. Sólo después de la carta dimensional esa
extensión se realiza como reloj cuántico finito de 108 pasos.

El levantamiento local de MD que porta la firma TRIT es la unidad
topológico--temporal de ese cristal genealógico.
Sobre
\(\mathcal H_{\rm clk}=\mathbb C^{108}\), con base
\(\{|j\rangle:j\in\mathbb Z/108\mathbb Z\}\), sean
\[
\zeta=e^{2\pi\ii/108},\qquad
|\widetilde k\rangle=\frac1{\sqrt{108}}
\sum_{j=0}^{107}\zeta^{jk}|j\rangle,
\qquad U_{\rm clk}|j\rangle=|j+1\rangle.
\]
\Needspace{12\baselineskip}
Con \(\omega_0=2\pi/(108t_0)\), el Hamiltoniano
\begin{equation}
H_{\rm clk}=\hbar_{\rm int}\omega_0
\sum_{k=0}^{107}k\,|\widetilde k\rangle\langle\widetilde k|
\label{p12-83-c3-s1:eq:trit-time-crystal-hamiltonian}
\end{equation}
genera exactamente
\begin{equation}
U_{\rm step}=\exp\!\left(
-\frac{\ii H_{\rm clk}t_0}{\hbar_{\rm int}}\right)=U_{\rm clk},
\qquad U_{\rm step}^{108}=I,
\label{p12-83-c3-s1:eq:trit-time-crystal-period}
\end{equation}
Si \(Z_{\rm clk}|j\rangle=\zeta^j|j\rangle\), entonces
\begin{equation}
\langle j_0|U_{\rm step}^{-n}Z_{\rm clk}U_{\rm step}^{n}|j_0\rangle
=\zeta^{j_0+n},
\label{p12-83-c3-s1:eq:trit-time-crystal-observable}
\end{equation}
cuya órbita posee periodo primitivo 108. En el vocabulario HMT, éste es el
cristal temporal genealógico interno: el calendario observable cierra,
\(C_{12}\) registra la vuelta y el estado TPK enriquecido no tiene por qué
retornar. La realización ulterior como fase física macroscópica exige,
además, protocolo de Floquet, respuesta subarmónica primitiva, robustez,
persistencia y un mecanismo de protección; esa carta no altera el teorema
finito de periodicidad unitaria.

\subsection{Publicación arquimediana, número genealógico y medición}

La publicación arquimediana utiliza cilindros ternarios semiabiertos
\begin{equation}
I_{N,L}=\left[\frac{N}{3^L},\frac{N+1}{3^L}\right),
\qquad \operatorname{diam}I_{N,L}=3^{-L}\longrightarrow0.
\label{p12-83-c3-s1:eq:ternary-cylinders}
\end{equation}
Una sección compatible determina cilindros encajados y, por completitud, un
único punto real. El número HMT es la sección misma,
\begin{equation}
x=(x_N)_N\in X_\chi
=\varprojlim_N\operatorname{Surv}^{(\chi)}_N,
\label{p12-83-c3-s1:eq:hmt-number-as-section}
\end{equation}
no la tupla formada por ella y sus evaluaciones. El valor aparece cuando un
lector olvida parte de la genealogía.
Por tanto, la precisión es profundidad: medir significa descender a un
cilindro más fino mientras el prefijo causal permanece en el antecedente.
La recta real se reconstruye como cociente arquimediano de un espacio de
historias mucho más rico.

\subsection{Realización dimensional del estado HMT: mapa de pantalla, magnitud observable y carta espectral}

La Mecánica Dimensional comienza donde el estado genealógico se acopla a una
frontera de realización. Fijado un cuerpo \(K_n\), sea \(M_n(x)\) el módulo
generado por los canales de prolongación de un estado \(x\), sea
\(\mathcal R_n(x)\) el submódulo de sus relaciones y sea \(\beta_n(x)\) el
dato de borde conservado. La pantalla dimensional es
\begin{equation}
\Sigma_n(x)=\bigl(K_n,M_n(x),\beta_n(x),\mathcal R_n(x)\bigr),
\qquad
d_{\Sigma_n}(x)=\dim_{K_n}\frac{M_n(x)}{\mathcal R_n(x)}.
\label{p12-83-c3-s1:eq:md-screen}
\end{equation}
Si \(M_n(x)\simeq K_n^{g_n}\) y las columnas de \(R_n(x)\) generan las
relaciones, entonces
\begin{equation}
\boxed{d_{\Sigma_n}(x)=g_n-\operatorname{rank}_{K_n}R_n(x).}
\label{p12-83-c3-s1:eq:md-rank}
\end{equation}
La dimensión no se adjunta como etiqueta: es el número de canales
independientes que sobreviven después de imponer las relaciones de frontera.

Una magnitud física es una sección de una recta dimensional \(L_D\); una
unidad es una carta \(\chi_u:L_D\to\mathbb R\) que publica su coordenada.
Por ello
\begin{equation}
\text{estado TPK}\longrightarrow\Sigma_n(x)
\longrightarrow s_D\in\Gamma(L_D)
\xrightarrow{\;\chi_u\;}[s_D]_u\in\mathbb R
\label{p12-83-c3-s1:eq:md-measurement-chain}
\end{equation}
separa de forma exacta estructura, tipo dimensional y número medido. Medir
consiste en acoplar una sección a un aparato reversible, imponer
compatibilidad de borde y aplicar el lector; cambiar de unidades cambia
\(\chi_u\), no la sección ni su genealogía.

\subsection{Doble proyección matemática y realización dimensional}

La rama excepcional conserva datos de calibre que la publicación real no
necesita. Sea \(\mathscr C_{\rm exc}\) el espacio de calibraciones y sea
\begin{equation}
\mathfrak E:
X_\infty^{\rm cal}\times\mathscr C_{\rm exc}
\longrightarrow\mathbf{Exc}
\label{p12-83-c3-s1:eq:exceptional-calibrated-reader}
\end{equation}
el lector de incidencia. Fijamos una vez
\begin{equation}
\begin{aligned}
\chi_{\rm exc}=\bigl(&\text{orden proyectivo},\ \text{calibre de Paley},
 \ \text{dodecada},\\
 &\text{alineamiento de incidencia},\ \text{cámara reticular}\bigr)
 \in\mathscr C_{\rm exc}
\end{aligned}
\label{p12-83-c3-s1:eq:exceptional-calibration}
\end{equation}
y abreviamos
\(\mathfrak E_{\chi_{\rm exc}}(x):=\mathfrak E(x,\chi_{\rm exc})\).
El calibre inicial se transporta por la historia; no vuelve a elegirse en
cada profundidad.

El mismo límite calibrado admite entonces dos lectores con codominios
distintos:
\begin{equation}
\mathbb R
\xleftarrow{\;\operatorname{ev}_{\mathbb R}\;}
X_\infty^{\rm cal}
\xrightarrow{\;\mathfrak E_{\chi_{\rm exc}}\;}
\mathbf{Exc}.
\label{p12-83-c3-s1:eq:double-projection}
\end{equation}
La primera rama contrae cilindros hasta un valor continuo. La segunda
conserva incidencias, orientaciones, calibre y frontera. Ambas descienden en
paralelo del mismo estado fuente. Esta bifurcación explica conjuntamente la
continuidad métrica y la emergencia de simetrías excepcionales. La segunda
rama es una reconstrucción calibrada: es exacta una vez fijado
\(\chi_{\rm exc}\), pero el estado desnudo no selecciona retrospectivamente
los cinco datos que componen ese calibre.

La realización dimensional no serializa esas dos ramas. Es un tercer lector
del mismo antecedente:
\begin{equation}
P_{\rm dim}:X_\infty\longrightarrow
\mathcal R_{\rm MD},
\qquad
\mathcal R_{\rm MD}=
\{\text{acción, masa, campo, geometría, información}\}.
\label{p12-83-c3-s1:eq:dimensional-publication}
\end{equation}
De este modo, un número, una simetría y una magnitud física pueden compartir
genealogía sin confundirse ni reducirse entre sí.

La tesis del continuo queda así formulada con precisión: la realidad
arquimediana es la sombra de radio cero de historias discretas infinitamente
refinables; aquello que la sombra omite reaparece como memoria, incidencia,
geometría y estructura física.

% RESULTADO_RECUPERADO. Owner íntegro de la dinámica del cociente, el retorno
% observable y los cociclos de modo y orientación. Se subordina a la única
% construcción APP--TRIT--TPK sin abrir un origen alternativo.
\begin{HMTScientificOwnerSubordinated}
\input{sections/hmt/05b_extensiones_dinamica_cociente}
\end{HMTScientificOwnerSubordinated}
}
